%% 中文简历模板 - [YOUR_NAME]
%% 适用于中国求职市场，A4 纸张，两页限制
%%
%% SETUP: 运行 /china setup 填写个人信息。
%% 编译命令：cd cv/chinese && lualatex -interaction=nonstopmode main_example.tex
%% 或：xelatex -interaction=nonstopmode main_example.tex
%%
%% 中文字体依赖：
%% - macOS: 建议使用 MacTeX，系统自带苹方/华文黑体
%% - Windows: 建议使用 MiKTeX/TeX Live，系统自带 SimHei/SimSun
%% - Linux: 建议安装 fonts-noto-cjk (Noto Sans CJK SC)
%% 如果编译报字体错误，请将下面 \setCJKmainfont 等改为系统已有的中文字体。

\documentclass[11pt,a4paper,sans]{moderncv}
\moderncvstyle{banking}
\moderncvcolor{blue}

% 强制姓名与章节标题使用 moderncv 强调色（蓝色）
\renewcommand*{\firstnamestyle}[1]{{\fontsize{34}{36}\bfseries\upshape\color{color1}#1}}
\renewcommand*{\lastnamestyle}[1]{{\fontsize{34}{36}\bfseries\upshape\color{color1}#1}}
\renewcommand*{\sectionstyle}[1]{{\sectionfont\color{color1}#1}}

\usepackage[utf8]{inputenc}
\usepackage{ctex}

% 如果系统能自动识别中文字体，可保留默认；否则取消下面注释并指定字体
% \setCJKmainfont{PingFang SC}       % macOS 示例
% \setCJKsansfont{PingFang SC}
% \setCJKmonofont{PingFang SC}

\usepackage{hyperref}
\hypersetup{
    colorlinks=true,
    linkcolor=blue,
    filecolor=magenta,
    urlcolor=blue,
    pdftitle={[YOUR_NAME] - 简历},
    pdfpagemode=FullScreen,
}
\usepackage[scale=0.77]{geometry}
\usepackage{import}
\usepackage{needspace}

% 个人信息
% 中文姓名通常放在姓名字段；如只有一个字段，可填入 \name{[姓名]}{}
\name{[姓名]}{}
\phone[mobile]{[手机号码]}
\email{[邮箱地址]}
\extrainfo{\href{[LinkedIn或GitHub链接]}{LinkedIn / GitHub} | 微信：[微信号]}

\begin{document}

\makecvtitle

% ============================================================
%     求职意向 / 个人简介
% ============================================================

\vspace{6pt}
\small{[3-5 行个人简介，针对目标岗位定制。说明你的核心优势、关键技能、以及能为公司带来的价值。示例：具有 X 年互联网产品研发经验的软件工程师，擅长 Python / Go 后端开发与分布式系统设计，熟悉云原生技术栈。主导过日活百万级产品的核心模块开发，具备良好的跨团队协作与技术决策能力。]}

% ============================================================
%     核心优势 / 技能
% ============================================================

\section{核心优势}
\vspace{1pt}
\begin{itemize}
\item \textbf{[技能类别 1]}: [具体技术栈/工具，如 Python、Go、Kubernetes、React 等]
\item \textbf{[技能类别 2]}: [具体能力，如分布式系统、机器学习、数据工程等]
\item \textbf{[技能类别 3]}: [领域专长，如金融科技、电商、SaaS 等]
\item \textbf{[技能类别 4]}: [软技能，如团队管理、跨部门协作、技术方案设计等]
\end{itemize}

% ============================================================
%     工作经历
% ============================================================

\section{工作经历}
\vspace{3pt}
\begin{itemize}

% --- 最近一份工作 ---
\item{\cventry{[YYYY--至今]}{[职位名称]}{[公司名称]}{[城市]}{}{
\vspace{1pt}
\begin{itemize}
    \item [职责或成果 1 - 尽量量化，如“负责 XX 系统，支撑日均 X 万订单”]
    \item [职责或成果 2]
    \item [职责或成果 3]
    \item [职责或成果 4]
\end{itemize}}}

\vspace{3pt}

% --- 上一份工作 ---
\item{\cventry{[YYYY--YYYY]}{[职位名称]}{[公司名称]}{[城市]}{}{
\vspace{1pt}
\begin{itemize}
    \item [职责或成果 1]
    \item [职责或成果 2]
    \item [职责或成果 3]
\end{itemize}}}

\vspace{3pt}

% --- 更早工作 ---
\item{\cventry{[YYYY--YYYY]}{[职位名称]}{[公司名称]}{[城市]}{}{
\vspace{1pt}
\begin{itemize}
    \item [职责或成果 1]
    \item [职责或成果 2]
\end{itemize}}}

\end{itemize}

% ============================================================
%     项目经历（可选）
% ============================================================

\section{项目经历}
\vspace{1pt}
\begin{itemize}
\item{\cventry{[YYYY--YYYY]}{[项目名称]}{[角色]}{[技术栈]}{}{
\vspace{1pt}
\begin{itemize}
    \item [项目背景与个人贡献]
    \item [量化成果]
\end{itemize}}}
\end{itemize}

% ============================================================
%     教育背景
% ============================================================

\section{教育背景}
\vspace{1pt}
\begin{itemize}
\item{\cventry{[YYYY--YYYY]}{[学历 / 专业]}{[学校名称]}{[城市]}{}{
\vspace{1pt}
[主修课程或荣誉，可选]
}}
\end{itemize}

% ============================================================
%     技能清单
% ============================================================

\section{技能清单}
\vspace{1pt}
\begin{itemize}
\item \textbf{编程语言}: [Python, Go, Java, TypeScript 等]
\item \textbf{框架/工具}: [Kubernetes, Docker, React, Django 等]
\item \textbf{语言}: [中文（母语）、英语（CET-6 / 流利）等]
\end{itemize}

% ============================================================
%     证书与荣誉（可选）
% ============================================================

\section{证书与荣誉}
\vspace{1pt}
\begin{itemize}
\item{\textbf{[证书/荣誉名称]} -- [颁发机构] ([年份])}
\end{itemize}

% ============================================================
%     自我评价（可选）
% ============================================================

\section{自我评价}
\vspace{1pt}
\begin{itemize}
\item [1-2 句话总结工作风格、职业素养或核心驱动力]
\end{itemize}

\end{document}
